\section*{Prompt: Merge Groups Across Batches}
\begin{description}
\item[Pipeline step:] Harmonization --- Hierarchical Merge
\item[Models:] Gemini 2.5 Pro
\item[Source:] \texttt{misc/prompts.py} $\rightarrow$ \texttt{MERGE\_GROUPS\_PROMPT}
\item[Called from:] \texttt{scripts/llm\_aggregator.py:\_merge\_groups\_batch()}
\end{description}
\begin{tcolorbox}[colback=yellow!8,colframe=yellow!50!black,left=4pt,right=4pt,top=4pt,bottom=4pt,arc=2pt]
\textbf{Note:} The \texttt{\{node\_type\}} and \texttt{\{output\_example\_json\}} placeholders are filled at runtime. The groups JSON is appended after this template.
\end{tcolorbox}

\begin{tcolorbox}[breakable,enhanced,colback=blue!3,colframe=blue!40!black,title={Prompt},fonttitle=\bfseries,left=8pt,right=8pt,top=6pt,bottom=6pt]

You are merging groups of \{node\_type\} that were created from different batches. Some groups may refer to the same concept even though they were processed separately.

Each input group has an \texttt{id}. Output only the merge mapping as \texttt{\{\{canonical\_name: [ids]\}\}}. All metadata (definitions, papers, snippets) is resolved from the IDs programmatically --- do NOT repeat them.

\textbf{Task:}
Review all groups below and merge any that refer to the same concept.

\textbf{Canonical name rules:} Use the most specific/common community-standard name. Sentence case. For measurement instruments, include the subset if applicable.

\textbf{Merge rules:}
\begin{itemize}
\item ONLY merge groups that refer to the EXACT SAME concept (true synonyms)
\item DO NOT merge hierarchical relationships --- "Reasoning" and "Mathematical Reasoning" are DIFFERENT and should stay separate
\item DO NOT merge parent-child relationships --- keep specialized variants separate
\item DO NOT merge distinct facets of a broader concept (e.g. "Mathematical reasoning" and "Commonsense reasoning" are separate facets, not synonyms)
\item For measurement instruments: DO NOT merge different datasets, subsets, or evaluation types (e.g. "HELM Math" $\neq$ "HELM Code"; "MMLU" $\neq$ "MMLU Math"; "User preference study" $\neq$ "Expert annotation evaluation")
\item Be VERY conservative: if unsure, keep them separate
\end{itemize}

\textbf{Output format:}
\begin{verbatim}
{{"groups": {{"Reasoning": [0, 3], "Mathematical reasoning": [1], "MMLU": [2, 5]}}, "ungrouped": [4, 6]}}
\end{verbatim}
\texttt{ungrouped} lists IDs that don't match any other group. Every input ID must appear exactly once.

\textbf{Input groups:}
\{groups\_json\}

Return only the JSON object.
\end{tcolorbox}
